\chapter{Épreuve chronométrée 5}\label{ch:iv1:timed-test-5}

\begin{ivtrains}
À faire d'une seule traite, sans notes ni calculatrice.
\end{ivtrains}

\ivtestsection{A}{6}{40}

\begin{question}\label{iq:iv1:timed-test-5:1}
Les événements arrivent selon un processus de Poisson de taux 2 par heure. Quelles sont l'espérance et la variance de l'instant d'arrivée du troisième événement~? Quelle est la probabilité que le troisième événement arrive après une heure~?
\end{question}

\begin{question}\label{iq:iv1:timed-test-5:2}
Vous lancez une pièce équilibrée jusqu'à obtenir le premier pile, et soit $N$ le nombre de lancers. Vous lancez ensuite $N$ dés équilibrés et additionnez les résultats. Déterminez l'espérance et la variance de la somme.
\end{question}

\begin{question}\label{iq:iv1:timed-test-5:3}
Une pièce a une probabilité inconnue de tomber sur pile, avec une loi a priori uniforme sur $[0, 1]$. Vous observez 3 piles en 4 lancers. Quelle est l'espérance a posteriori~? Quelle est la probabilité que les deux lancers suivants donnent tous deux pile~?
\end{question}

\begin{question}\label{iq:iv1:timed-test-5:4}
Cinq boules rouges et cinq boules bleues sont alignées dans un ordre aléatoire. Une série est un bloc maximal d'une même couleur. Quelle est l'espérance du nombre de séries~? Quelle est la probabilité qu'il y ait exactement deux séries~?
\end{question}

\begin{question}\label{iq:iv1:timed-test-5:5}
On lance trois dés équilibrés. Quelle est la probabilité que le produit des trois nombres obtenus soit divisible par 6~?
\end{question}

\begin{question}\label{iq:iv1:timed-test-5:6}
$X$ et $Y$ sont des lois normales centrées réduites indépendantes. Quelle est la probabilité que le point $(X, Y)$ se trouve à une distance au plus égale à 1 de l'origine~? Quel rayon contient exactement la moitié de la probabilité~?
\end{question}

\ivtestsection{B}{6}{60}

\begin{question}\label{iq:iv1:timed-test-5:7}
Une élection opposant A et B se termine par 7 voix pour A et 3 pour B, dépouillées dans un ordre aléatoire. Quelle est la probabilité que A soit strictement en tête tout au long du dépouillement~?
\end{question}

\begin{question}\label{iq:iv1:timed-test-5:8}
Dans un graphe aléatoire à 10 sommets, chaque arête est présente indépendamment avec probabilité $\frac12$. Déterminez l'espérance et la variance du nombre de triangles.
\end{question}

\begin{question}\label{iq:iv1:timed-test-5:9}
En utilisant $\log_{10}2 = 0.30103$, combien de chiffres décimaux possède $2^{1000}$, et quel est son premier chiffre~?
\end{question}

\begin{question}\label{iq:iv1:timed-test-5:10}
La ligne de bus A passe toutes les 10 minutes et la ligne B toutes les 15 minutes, chacune selon un horaire fixe, avec des phases uniformes indépendantes. Vous arrivez à un instant aléatoire et prenez le premier bus. Quel est votre temps d'attente espéré~?
\end{question}

\begin{question}\label{iq:iv1:timed-test-5:11}
Un jeu télévisé comporte 5 portes, dont une cache un lot. Vous choisissez une porte. L'animateur, qui sait où se trouve le lot, ouvre une autre porte derrière laquelle il n'y a pas de lot et vous propose de changer de porte. Si vous changez au hasard pour l'une des trois portes restantes fermées, quelle est votre probabilité de gagner~?
\end{question}

\begin{question}\label{iq:iv1:timed-test-5:12}
$X$ suit une loi binomiale à 100 essais, de probabilité de succès $\frac12$. Majorez $P(X \ge 70)$ par l'inégalité de Hoeffding $\exp(-2t^2/n)$, où $t$ désigne l'écart à la moyenne. Comparez cette majoration à l'inégalité de Tchebychev et à l'approximation normale.
\end{question}
